% Options for packages loaded elsewhere
% Options for packages loaded elsewhere
\PassOptionsToPackage{unicode}{hyperref}
\PassOptionsToPackage{hyphens}{url}
\PassOptionsToPackage{dvipsnames,svgnames,x11names}{xcolor}
%
\documentclass[
  letterpaper,
  DIV=11,
  numbers=noendperiod]{scrreprt}
\usepackage{xcolor}
\usepackage{amsmath,amssymb}
\setcounter{secnumdepth}{5}
\usepackage{iftex}
\ifPDFTeX
  \usepackage[T1]{fontenc}
  \usepackage[utf8]{inputenc}
  \usepackage{textcomp} % provide euro and other symbols
\else % if luatex or xetex
  \usepackage{unicode-math} % this also loads fontspec
  \defaultfontfeatures{Scale=MatchLowercase}
  \defaultfontfeatures[\rmfamily]{Ligatures=TeX,Scale=1}
\fi
\usepackage{lmodern}
\ifPDFTeX\else
  % xetex/luatex font selection
\fi
% Use upquote if available, for straight quotes in verbatim environments
\IfFileExists{upquote.sty}{\usepackage{upquote}}{}
\IfFileExists{microtype.sty}{% use microtype if available
  \usepackage[]{microtype}
  \UseMicrotypeSet[protrusion]{basicmath} % disable protrusion for tt fonts
}{}
\makeatletter
\@ifundefined{KOMAClassName}{% if non-KOMA class
  \IfFileExists{parskip.sty}{%
    \usepackage{parskip}
  }{% else
    \setlength{\parindent}{0pt}
    \setlength{\parskip}{6pt plus 2pt minus 1pt}}
}{% if KOMA class
  \KOMAoptions{parskip=half}}
\makeatother
% Make \paragraph and \subparagraph free-standing
\makeatletter
\ifx\paragraph\undefined\else
  \let\oldparagraph\paragraph
  \renewcommand{\paragraph}{
    \@ifstar
      \xxxParagraphStar
      \xxxParagraphNoStar
  }
  \newcommand{\xxxParagraphStar}[1]{\oldparagraph*{#1}\mbox{}}
  \newcommand{\xxxParagraphNoStar}[1]{\oldparagraph{#1}\mbox{}}
\fi
\ifx\subparagraph\undefined\else
  \let\oldsubparagraph\subparagraph
  \renewcommand{\subparagraph}{
    \@ifstar
      \xxxSubParagraphStar
      \xxxSubParagraphNoStar
  }
  \newcommand{\xxxSubParagraphStar}[1]{\oldsubparagraph*{#1}\mbox{}}
  \newcommand{\xxxSubParagraphNoStar}[1]{\oldsubparagraph{#1}\mbox{}}
\fi
\makeatother

\usepackage{color}
\usepackage{fancyvrb}
\newcommand{\VerbBar}{|}
\newcommand{\VERB}{\Verb[commandchars=\\\{\}]}
\DefineVerbatimEnvironment{Highlighting}{Verbatim}{commandchars=\\\{\}}
% Add ',fontsize=\small' for more characters per line
\usepackage{framed}
\definecolor{shadecolor}{RGB}{241,243,245}
\newenvironment{Shaded}{\begin{snugshade}}{\end{snugshade}}
\newcommand{\AlertTok}[1]{\textcolor[rgb]{0.68,0.00,0.00}{#1}}
\newcommand{\AnnotationTok}[1]{\textcolor[rgb]{0.37,0.37,0.37}{#1}}
\newcommand{\AttributeTok}[1]{\textcolor[rgb]{0.40,0.45,0.13}{#1}}
\newcommand{\BaseNTok}[1]{\textcolor[rgb]{0.68,0.00,0.00}{#1}}
\newcommand{\BuiltInTok}[1]{\textcolor[rgb]{0.00,0.23,0.31}{#1}}
\newcommand{\CharTok}[1]{\textcolor[rgb]{0.13,0.47,0.30}{#1}}
\newcommand{\CommentTok}[1]{\textcolor[rgb]{0.37,0.37,0.37}{#1}}
\newcommand{\CommentVarTok}[1]{\textcolor[rgb]{0.37,0.37,0.37}{\textit{#1}}}
\newcommand{\ConstantTok}[1]{\textcolor[rgb]{0.56,0.35,0.01}{#1}}
\newcommand{\ControlFlowTok}[1]{\textcolor[rgb]{0.00,0.23,0.31}{\textbf{#1}}}
\newcommand{\DataTypeTok}[1]{\textcolor[rgb]{0.68,0.00,0.00}{#1}}
\newcommand{\DecValTok}[1]{\textcolor[rgb]{0.68,0.00,0.00}{#1}}
\newcommand{\DocumentationTok}[1]{\textcolor[rgb]{0.37,0.37,0.37}{\textit{#1}}}
\newcommand{\ErrorTok}[1]{\textcolor[rgb]{0.68,0.00,0.00}{#1}}
\newcommand{\ExtensionTok}[1]{\textcolor[rgb]{0.00,0.23,0.31}{#1}}
\newcommand{\FloatTok}[1]{\textcolor[rgb]{0.68,0.00,0.00}{#1}}
\newcommand{\FunctionTok}[1]{\textcolor[rgb]{0.28,0.35,0.67}{#1}}
\newcommand{\ImportTok}[1]{\textcolor[rgb]{0.00,0.46,0.62}{#1}}
\newcommand{\InformationTok}[1]{\textcolor[rgb]{0.37,0.37,0.37}{#1}}
\newcommand{\KeywordTok}[1]{\textcolor[rgb]{0.00,0.23,0.31}{\textbf{#1}}}
\newcommand{\NormalTok}[1]{\textcolor[rgb]{0.00,0.23,0.31}{#1}}
\newcommand{\OperatorTok}[1]{\textcolor[rgb]{0.37,0.37,0.37}{#1}}
\newcommand{\OtherTok}[1]{\textcolor[rgb]{0.00,0.23,0.31}{#1}}
\newcommand{\PreprocessorTok}[1]{\textcolor[rgb]{0.68,0.00,0.00}{#1}}
\newcommand{\RegionMarkerTok}[1]{\textcolor[rgb]{0.00,0.23,0.31}{#1}}
\newcommand{\SpecialCharTok}[1]{\textcolor[rgb]{0.37,0.37,0.37}{#1}}
\newcommand{\SpecialStringTok}[1]{\textcolor[rgb]{0.13,0.47,0.30}{#1}}
\newcommand{\StringTok}[1]{\textcolor[rgb]{0.13,0.47,0.30}{#1}}
\newcommand{\VariableTok}[1]{\textcolor[rgb]{0.07,0.07,0.07}{#1}}
\newcommand{\VerbatimStringTok}[1]{\textcolor[rgb]{0.13,0.47,0.30}{#1}}
\newcommand{\WarningTok}[1]{\textcolor[rgb]{0.37,0.37,0.37}{\textit{#1}}}

\usepackage{longtable,booktabs,array}
\newcounter{none} % for unnumbered tables
\usepackage{calc} % for calculating minipage widths
% Correct order of tables after \paragraph or \subparagraph
\usepackage{etoolbox}
\makeatletter
\patchcmd\longtable{\par}{\if@noskipsec\mbox{}\fi\par}{}{}
\makeatother
% Allow footnotes in longtable head/foot
\IfFileExists{footnotehyper.sty}{\usepackage{footnotehyper}}{\usepackage{footnote}}
\makesavenoteenv{longtable}
\usepackage{graphicx}
\makeatletter
\newsavebox\pandoc@box
\newcommand*\pandocbounded[1]{% scales image to fit in text height/width
  \sbox\pandoc@box{#1}%
  \Gscale@div\@tempa{\textheight}{\dimexpr\ht\pandoc@box+\dp\pandoc@box\relax}%
  \Gscale@div\@tempb{\linewidth}{\wd\pandoc@box}%
  \ifdim\@tempb\p@<\@tempa\p@\let\@tempa\@tempb\fi% select the smaller of both
  \ifdim\@tempa\p@<\p@\scalebox{\@tempa}{\usebox\pandoc@box}%
  \else\usebox{\pandoc@box}%
  \fi%
}
% Set default figure placement to htbp
\def\fps@figure{htbp}
\makeatother





\setlength{\emergencystretch}{3em} % prevent overfull lines

\providecommand{\tightlist}{%
  \setlength{\itemsep}{0pt}\setlength{\parskip}{0pt}}



 


\KOMAoption{captions}{tableheading}
\makeatletter
\@ifpackageloaded{bookmark}{}{\usepackage{bookmark}}
\makeatother
\makeatletter
\@ifpackageloaded{caption}{}{\usepackage{caption}}
\AtBeginDocument{%
\ifdefined\contentsname
  \renewcommand*\contentsname{Table of contents}
\else
  \newcommand\contentsname{Table of contents}
\fi
\ifdefined\listfigurename
  \renewcommand*\listfigurename{List of Figures}
\else
  \newcommand\listfigurename{List of Figures}
\fi
\ifdefined\listtablename
  \renewcommand*\listtablename{List of Tables}
\else
  \newcommand\listtablename{List of Tables}
\fi
\ifdefined\figurename
  \renewcommand*\figurename{Figure}
\else
  \newcommand\figurename{Figure}
\fi
\ifdefined\tablename
  \renewcommand*\tablename{Table}
\else
  \newcommand\tablename{Table}
\fi
}
\@ifpackageloaded{float}{}{\usepackage{float}}
\floatstyle{ruled}
\@ifundefined{c@chapter}{\newfloat{codelisting}{h}{lop}}{\newfloat{codelisting}{h}{lop}[chapter]}
\floatname{codelisting}{Listing}
\newcommand*\listoflistings{\listof{codelisting}{List of Listings}}
\makeatother
\makeatletter
\makeatother
\makeatletter
\@ifpackageloaded{caption}{}{\usepackage{caption}}
\@ifpackageloaded{subcaption}{}{\usepackage{subcaption}}
\makeatother
\usepackage{bookmark}
\IfFileExists{xurl.sty}{\usepackage{xurl}}{} % add URL line breaks if available
\urlstyle{same}
\hypersetup{
  pdftitle={Analisis del Cambio Climático con R},
  pdfauthor={Junior Ponce},
  colorlinks=true,
  linkcolor={blue},
  filecolor={Maroon},
  citecolor={Blue},
  urlcolor={Blue},
  pdfcreator={LaTeX via pandoc}}


\title{Analisis del Cambio Climático con R}
\author{Junior Ponce}
\date{2026-05-29}
\begin{document}
\maketitle

\renewcommand*\contentsname{Table of contents}
{
\setcounter{tocdepth}{2}
\tableofcontents
}

\bookmarksetup{startatroot}

\chapter*{Bienvenida}\label{bienvenida}
\addcontentsline{toc}{chapter}{Bienvenida}

\markboth{Bienvenida}{Bienvenida}

\bookmarksetup{startatroot}

\chapter{Capítulo 1: Datos de temperatura global
(NASA)}\label{capuxedtulo-1-datos-de-temperatura-global-nasa}

En este capítulo vamos a descargar y explorar el dataset
\textbf{GISTEMP} de la NASA, que registra las anomalías de temperatura
global desde 1880.

\begin{quote}
Una \textbf{anomalía de temperatura} es la diferencia entre la
temperatura de un año y el promedio histórico (1951--1980). Si es
positiva, ese año fue más cálido de lo normal.
\end{quote}

\section{Cargar paquetes}\label{cargar-paquetes}

\begin{Shaded}
\begin{Highlighting}[]
\FunctionTok{library}\NormalTok{(tidyverse)}
\FunctionTok{library}\NormalTok{(lubridate)}
\FunctionTok{library}\NormalTok{(janitor)}
\FunctionTok{library}\NormalTok{(skimr)}
\end{Highlighting}
\end{Shaded}

\section{Descargar los datos}\label{descargar-los-datos}

\begin{Shaded}
\begin{Highlighting}[]
\NormalTok{url }\OtherTok{\textless{}{-}} \StringTok{"https://data.giss.nasa.gov/gistemp/tabledata\_v4/GLB.Ts+dSST.csv"}

\NormalTok{clima\_raw }\OtherTok{\textless{}{-}} \FunctionTok{read\_csv}\NormalTok{(}\StringTok{"GLB.Ts+dSST.csv"}\NormalTok{, }\AttributeTok{skip =} \DecValTok{1}\NormalTok{, }\AttributeTok{na =} \StringTok{"***"}\NormalTok{)}

\FunctionTok{glimpse}\NormalTok{(clima\_raw)}
\end{Highlighting}
\end{Shaded}

\begin{verbatim}
Rows: 147
Columns: 19
$ Year  <dbl> 1880, 1881, 1882, 1883, 1884, 1885, 1886, 1887, 1888, 1889, 1890~
$ Jan   <dbl> -0.19, -0.20, 0.16, -0.29, -0.13, -0.59, -0.44, -0.72, -0.35, -0~
$ Feb   <dbl> -0.25, -0.16, 0.14, -0.36, -0.08, -0.34, -0.51, -0.57, -0.37, 0.~
$ Mar   <dbl> -0.10, 0.02, 0.04, -0.12, -0.36, -0.27, -0.43, -0.36, -0.42, 0.0~
$ Apr   <dbl> -0.17, 0.04, -0.16, -0.18, -0.40, -0.42, -0.28, -0.35, -0.21, 0.~
$ May   <dbl> -0.10, 0.06, -0.14, -0.17, -0.34, -0.46, -0.25, -0.31, -0.22, -0~
$ Jun   <dbl> -0.21, -0.19, -0.22, -0.07, -0.35, -0.44, -0.35, -0.25, -0.18, -~
$ Jul   <dbl> -0.19, 0.01, -0.16, -0.07, -0.30, -0.34, -0.19, -0.26, -0.11, -0~
$ Aug   <dbl> -0.11, -0.04, -0.08, -0.14, -0.28, -0.32, -0.31, -0.36, -0.16, -~
$ Sep   <dbl> -0.15, -0.15, -0.15, -0.22, -0.27, -0.28, -0.24, -0.26, -0.12, -~
$ Oct   <dbl> -0.24, -0.22, -0.23, -0.11, -0.25, -0.23, -0.28, -0.36, 0.01, -0~
$ Nov   <dbl> -0.22, -0.19, -0.17, -0.24, -0.34, -0.24, -0.28, -0.26, 0.02, -0~
$ Dec   <dbl> -0.19, -0.07, -0.36, -0.11, -0.31, -0.10, -0.26, -0.33, -0.05, -~
$ `J-D` <dbl> -0.18, -0.09, -0.11, -0.17, -0.28, -0.34, -0.32, -0.37, -0.18, -~
$ `D-N` <dbl> NA, -0.10, -0.09, -0.20, -0.27, -0.35, -0.31, -0.36, -0.20, -0.0~
$ DJF   <dbl> NA, -0.18, 0.08, -0.34, -0.11, -0.41, -0.35, -0.52, -0.35, 0.00,~
$ MAM   <dbl> -0.12, 0.04, -0.08, -0.16, -0.36, -0.38, -0.32, -0.34, -0.28, 0.~
$ JJA   <dbl> -0.17, -0.07, -0.15, -0.09, -0.31, -0.37, -0.28, -0.29, -0.15, -~
$ SON   <dbl> -0.20, -0.18, -0.18, -0.19, -0.28, -0.25, -0.27, -0.29, -0.03, -~
\end{verbatim}

\section{Limpiar los datos}\label{limpiar-los-datos}

\begin{Shaded}
\begin{Highlighting}[]
\NormalTok{clima }\OtherTok{\textless{}{-}}\NormalTok{ clima\_raw }\SpecialCharTok{|\textgreater{}}
  \FunctionTok{clean\_names}\NormalTok{() }\SpecialCharTok{|\textgreater{}}
  \FunctionTok{select}\NormalTok{(year, jan}\SpecialCharTok{:}\NormalTok{dec) }\SpecialCharTok{|\textgreater{}}
  \FunctionTok{pivot\_longer}\NormalTok{(}\AttributeTok{cols =}\NormalTok{ jan}\SpecialCharTok{:}\NormalTok{dec,}
               \AttributeTok{names\_to  =} \StringTok{"mes"}\NormalTok{,}
               \AttributeTok{values\_to =} \StringTok{"anomalia"}\NormalTok{) }\SpecialCharTok{|\textgreater{}}
  \FunctionTok{mutate}\NormalTok{(}
    \AttributeTok{mes   =} \FunctionTok{match}\NormalTok{(mes, }\FunctionTok{c}\NormalTok{(}\StringTok{"jan"}\NormalTok{,}\StringTok{"feb"}\NormalTok{,}\StringTok{"mar"}\NormalTok{,}\StringTok{"apr"}\NormalTok{,}
                         \StringTok{"may"}\NormalTok{,}\StringTok{"jun"}\NormalTok{,}\StringTok{"jul"}\NormalTok{,}\StringTok{"aug"}\NormalTok{,}
                         \StringTok{"sep"}\NormalTok{,}\StringTok{"oct"}\NormalTok{,}\StringTok{"nov"}\NormalTok{,}\StringTok{"dec"}\NormalTok{)),}
    \AttributeTok{fecha =} \FunctionTok{make\_date}\NormalTok{(year, mes, }\DecValTok{1}\NormalTok{)}
\NormalTok{  ) }\SpecialCharTok{|\textgreater{}}
  \FunctionTok{drop\_na}\NormalTok{(anomalia)}

\FunctionTok{glimpse}\NormalTok{(clima)}
\end{Highlighting}
\end{Shaded}

\begin{verbatim}
Rows: 1,756
Columns: 4
$ year     <dbl> 1880, 1880, 1880, 1880, 1880, 1880, 1880, 1880, 1880, 1880, 1~
$ mes      <int> 1, 2, 3, 4, 5, 6, 7, 8, 9, 10, 11, 12, 1, 2, 3, 4, 5, 6, 7, 8~
$ anomalia <dbl> -0.19, -0.25, -0.10, -0.17, -0.10, -0.21, -0.19, -0.11, -0.15~
$ fecha    <date> 1880-01-01, 1880-02-01, 1880-03-01, 1880-04-01, 1880-05-01, ~
\end{verbatim}

\section{Primer vistazo}\label{primer-vistazo}

\begin{Shaded}
\begin{Highlighting}[]
\FunctionTok{skim}\NormalTok{(clima)}
\end{Highlighting}
\end{Shaded}

\begin{longtable}[]{@{}ll@{}}
\caption{Data summary}\tabularnewline
\toprule\noalign{}
\endfirsthead
\endhead
\bottomrule\noalign{}
\endlastfoot
Name & clima \\
Number of rows & 1756 \\
Number of columns & 4 \\
\_\_\_\_\_\_\_\_\_\_\_\_\_\_\_\_\_\_\_\_\_\_\_ & \\
Column type frequency: & \\
Date & 1 \\
numeric & 3 \\
\_\_\_\_\_\_\_\_\_\_\_\_\_\_\_\_\_\_\_\_\_\_\_\_ & \\
Group variables & None \\
\end{longtable}

\textbf{Variable type: Date}

{\def\LTcaptype{none} % do not increment counter
\begin{longtable}[]{@{}
  >{\raggedright\arraybackslash}p{(\linewidth - 12\tabcolsep) * \real{0.1750}}
  >{\raggedleft\arraybackslash}p{(\linewidth - 12\tabcolsep) * \real{0.1250}}
  >{\raggedleft\arraybackslash}p{(\linewidth - 12\tabcolsep) * \real{0.1750}}
  >{\raggedright\arraybackslash}p{(\linewidth - 12\tabcolsep) * \real{0.1375}}
  >{\raggedright\arraybackslash}p{(\linewidth - 12\tabcolsep) * \real{0.1375}}
  >{\raggedright\arraybackslash}p{(\linewidth - 12\tabcolsep) * \real{0.1375}}
  >{\raggedleft\arraybackslash}p{(\linewidth - 12\tabcolsep) * \real{0.1125}}@{}}
\toprule\noalign{}
\begin{minipage}[b]{\linewidth}\raggedright
skim\_variable
\end{minipage} & \begin{minipage}[b]{\linewidth}\raggedleft
n\_missing
\end{minipage} & \begin{minipage}[b]{\linewidth}\raggedleft
complete\_rate
\end{minipage} & \begin{minipage}[b]{\linewidth}\raggedright
min
\end{minipage} & \begin{minipage}[b]{\linewidth}\raggedright
max
\end{minipage} & \begin{minipage}[b]{\linewidth}\raggedright
median
\end{minipage} & \begin{minipage}[b]{\linewidth}\raggedleft
n\_unique
\end{minipage} \\
\midrule\noalign{}
\endhead
\bottomrule\noalign{}
\endlastfoot
fecha & 0 & 1 & 1880-01-01 & 2026-04-01 & 1953-02-15 & 1756 \\
\end{longtable}
}

\textbf{Variable type: numeric}

{\def\LTcaptype{none} % do not increment counter
\begin{longtable}[]{@{}
  >{\raggedright\arraybackslash}p{(\linewidth - 20\tabcolsep) * \real{0.1429}}
  >{\raggedleft\arraybackslash}p{(\linewidth - 20\tabcolsep) * \real{0.1020}}
  >{\raggedleft\arraybackslash}p{(\linewidth - 20\tabcolsep) * \real{0.1429}}
  >{\raggedleft\arraybackslash}p{(\linewidth - 20\tabcolsep) * \real{0.0816}}
  >{\raggedleft\arraybackslash}p{(\linewidth - 20\tabcolsep) * \real{0.0612}}
  >{\raggedleft\arraybackslash}p{(\linewidth - 20\tabcolsep) * \real{0.0816}}
  >{\raggedleft\arraybackslash}p{(\linewidth - 20\tabcolsep) * \real{0.0816}}
  >{\raggedleft\arraybackslash}p{(\linewidth - 20\tabcolsep) * \real{0.0816}}
  >{\raggedleft\arraybackslash}p{(\linewidth - 20\tabcolsep) * \real{0.0816}}
  >{\raggedleft\arraybackslash}p{(\linewidth - 20\tabcolsep) * \real{0.0816}}
  >{\raggedright\arraybackslash}p{(\linewidth - 20\tabcolsep) * \real{0.0612}}@{}}
\toprule\noalign{}
\begin{minipage}[b]{\linewidth}\raggedright
skim\_variable
\end{minipage} & \begin{minipage}[b]{\linewidth}\raggedleft
n\_missing
\end{minipage} & \begin{minipage}[b]{\linewidth}\raggedleft
complete\_rate
\end{minipage} & \begin{minipage}[b]{\linewidth}\raggedleft
mean
\end{minipage} & \begin{minipage}[b]{\linewidth}\raggedleft
sd
\end{minipage} & \begin{minipage}[b]{\linewidth}\raggedleft
p0
\end{minipage} & \begin{minipage}[b]{\linewidth}\raggedleft
p25
\end{minipage} & \begin{minipage}[b]{\linewidth}\raggedleft
p50
\end{minipage} & \begin{minipage}[b]{\linewidth}\raggedleft
p75
\end{minipage} & \begin{minipage}[b]{\linewidth}\raggedleft
p100
\end{minipage} & \begin{minipage}[b]{\linewidth}\raggedright
hist
\end{minipage} \\
\midrule\noalign{}
\endhead
\bottomrule\noalign{}
\endlastfoot
year & 0 & 1 & 1952.67 & 42.26 & 1880.00 & 1916.00 & 1953.00 & 1989.00 &
2026.00 & ▇▇▇▇▇ \\
mes & 0 & 1 & 6.49 & 3.45 & 1.00 & 3.00 & 6.00 & 9.00 & 12.00 & ▇▅▅▅▇ \\
anomalia & 0 & 1 & 0.08 & 0.42 & -0.82 & -0.22 & -0.03 & 0.32 & 1.48 &
▂▇▃▂▁ \\
\end{longtable}
}

\bookmarksetup{startatroot}

\chapter{Capítulo 2: Visualizaciones con
ggplot2}\label{capuxedtulo-2-visualizaciones-con-ggplot2}

Ahora que tenemos los datos limpios, vamos a construir visualizaciones
que muestren cómo ha cambiado la temperatura global desde 1880.

\section{Cargar paquetes y datos}\label{cargar-paquetes-y-datos}

\begin{Shaded}
\begin{Highlighting}[]
\FunctionTok{library}\NormalTok{(tidyverse)}
\FunctionTok{library}\NormalTok{(lubridate)}
\FunctionTok{library}\NormalTok{(janitor)}
\FunctionTok{library}\NormalTok{(scales)}
\FunctionTok{library}\NormalTok{(patchwork)}
\FunctionTok{library}\NormalTok{(ggthemes)}

\NormalTok{clima\_raw }\OtherTok{\textless{}{-}} \FunctionTok{read\_csv}\NormalTok{(}\StringTok{"GLB.Ts+dSST.csv"}\NormalTok{, }\AttributeTok{skip =} \DecValTok{1}\NormalTok{, }\AttributeTok{na =} \StringTok{"***"}\NormalTok{)}

\NormalTok{clima }\OtherTok{\textless{}{-}}\NormalTok{ clima\_raw }\SpecialCharTok{|\textgreater{}}
  \FunctionTok{clean\_names}\NormalTok{() }\SpecialCharTok{|\textgreater{}}
  \FunctionTok{select}\NormalTok{(year, jan}\SpecialCharTok{:}\NormalTok{dec) }\SpecialCharTok{|\textgreater{}}
  \FunctionTok{pivot\_longer}\NormalTok{(}\AttributeTok{cols =}\NormalTok{ jan}\SpecialCharTok{:}\NormalTok{dec,}
               \AttributeTok{names\_to  =} \StringTok{"mes"}\NormalTok{,}
               \AttributeTok{values\_to =} \StringTok{"anomalia"}\NormalTok{) }\SpecialCharTok{|\textgreater{}}
  \FunctionTok{mutate}\NormalTok{(}
    \AttributeTok{mes   =} \FunctionTok{match}\NormalTok{(mes, }\FunctionTok{c}\NormalTok{(}\StringTok{"jan"}\NormalTok{,}\StringTok{"feb"}\NormalTok{,}\StringTok{"mar"}\NormalTok{,}\StringTok{"apr"}\NormalTok{,}
                         \StringTok{"may"}\NormalTok{,}\StringTok{"jun"}\NormalTok{,}\StringTok{"jul"}\NormalTok{,}\StringTok{"aug"}\NormalTok{,}
                         \StringTok{"sep"}\NormalTok{,}\StringTok{"oct"}\NormalTok{,}\StringTok{"nov"}\NormalTok{,}\StringTok{"dec"}\NormalTok{)),}
    \AttributeTok{fecha =} \FunctionTok{make\_date}\NormalTok{(year, mes, }\DecValTok{1}\NormalTok{)}
\NormalTok{  ) }\SpecialCharTok{|\textgreater{}}
  \FunctionTok{drop\_na}\NormalTok{(anomalia)}
\end{Highlighting}
\end{Shaded}

\section{Gráfico 1 --- Serie de tiempo
completa}\label{gruxe1fico-1-serie-de-tiempo-completa}

\begin{Shaded}
\begin{Highlighting}[]
\NormalTok{clima }\SpecialCharTok{|\textgreater{}}
  \FunctionTok{group\_by}\NormalTok{(year) }\SpecialCharTok{|\textgreater{}}
  \FunctionTok{summarise}\NormalTok{(}\AttributeTok{anomalia\_anual =} \FunctionTok{mean}\NormalTok{(anomalia)) }\SpecialCharTok{|\textgreater{}}
  \FunctionTok{ggplot}\NormalTok{(}\FunctionTok{aes}\NormalTok{(}\AttributeTok{x =}\NormalTok{ year, }\AttributeTok{y =}\NormalTok{ anomalia\_anual)) }\SpecialCharTok{+}
  \FunctionTok{geom\_line}\NormalTok{(}\AttributeTok{color =} \StringTok{"\#2196F3"}\NormalTok{, }\AttributeTok{linewidth =} \FloatTok{0.7}\NormalTok{) }\SpecialCharTok{+}
  \FunctionTok{geom\_smooth}\NormalTok{(}\AttributeTok{method =} \StringTok{"loess"}\NormalTok{, }\AttributeTok{color =} \StringTok{"\#E53935"}\NormalTok{,}
              \AttributeTok{se =} \ConstantTok{FALSE}\NormalTok{, }\AttributeTok{linewidth =} \FloatTok{1.2}\NormalTok{) }\SpecialCharTok{+}
  \FunctionTok{geom\_hline}\NormalTok{(}\AttributeTok{yintercept =} \DecValTok{0}\NormalTok{, }\AttributeTok{linetype =} \StringTok{"dashed"}\NormalTok{,}
             \AttributeTok{color =} \StringTok{"gray50"}\NormalTok{) }\SpecialCharTok{+}
  \FunctionTok{labs}\NormalTok{(}
    \AttributeTok{title    =} \StringTok{"Anomalía de temperatura global (1880–2024)"}\NormalTok{,}
    \AttributeTok{subtitle =} \StringTok{"Línea azul: dato anual · Línea roja: tendencia"}\NormalTok{,}
    \AttributeTok{x        =} \StringTok{"Año"}\NormalTok{,}
    \AttributeTok{y        =} \StringTok{"Anomalía (°C)"}\NormalTok{,}
    \AttributeTok{caption  =} \StringTok{"Fuente: NASA GISS Surface Temperature Analysis"}
\NormalTok{  ) }\SpecialCharTok{+}
  \FunctionTok{theme\_minimal}\NormalTok{()}
\end{Highlighting}
\end{Shaded}

\begin{verbatim}
`geom_smooth()` using formula = 'y ~ x'
\end{verbatim}

\begin{figure}[H]

{\centering \pandocbounded{\includegraphics[keepaspectratio]{visualizaciones_files/figure-pdf/unnamed-chunk-2-1.pdf}}

}

\caption{Anomalías de temperatura global 1880--2024}

\end{figure}%

\section{Gráfico 2 --- Mapa de calor por mes y
año}\label{gruxe1fico-2-mapa-de-calor-por-mes-y-auxf1o}

\begin{Shaded}
\begin{Highlighting}[]
\NormalTok{clima }\SpecialCharTok{|\textgreater{}}
  \FunctionTok{filter}\NormalTok{(year }\SpecialCharTok{\textgreater{}=} \DecValTok{1950}\NormalTok{) }\SpecialCharTok{|\textgreater{}}
  \FunctionTok{ggplot}\NormalTok{(}\FunctionTok{aes}\NormalTok{(}\AttributeTok{x =}\NormalTok{ year, }\AttributeTok{y =} \FunctionTok{factor}\NormalTok{(mes), }\AttributeTok{fill =}\NormalTok{ anomalia)) }\SpecialCharTok{+}
  \FunctionTok{geom\_tile}\NormalTok{() }\SpecialCharTok{+}
  \FunctionTok{scale\_fill\_gradient2}\NormalTok{(}
    \AttributeTok{low      =} \StringTok{"\#2166AC"}\NormalTok{,}
    \AttributeTok{mid      =} \StringTok{"white"}\NormalTok{,}
    \AttributeTok{high     =} \StringTok{"\#D6604D"}\NormalTok{,}
    \AttributeTok{midpoint =} \DecValTok{0}
\NormalTok{  ) }\SpecialCharTok{+}
  \FunctionTok{scale\_y\_discrete}\NormalTok{(}\AttributeTok{labels =} \FunctionTok{c}\NormalTok{(}\StringTok{"Ene"}\NormalTok{,}\StringTok{"Feb"}\NormalTok{,}\StringTok{"Mar"}\NormalTok{,}\StringTok{"Abr"}\NormalTok{,}\StringTok{"May"}\NormalTok{,}
                               \StringTok{"Jun"}\NormalTok{,}\StringTok{"Jul"}\NormalTok{,}\StringTok{"Ago"}\NormalTok{,}\StringTok{"Sep"}\NormalTok{,}\StringTok{"Oct"}\NormalTok{,}
                               \StringTok{"Nov"}\NormalTok{,}\StringTok{"Dic"}\NormalTok{)) }\SpecialCharTok{+}
  \FunctionTok{labs}\NormalTok{(}
    \AttributeTok{title   =} \StringTok{"Anomalías de temperatura por mes (1950–2024)"}\NormalTok{,}
    \AttributeTok{x       =} \StringTok{"Año"}\NormalTok{,}
    \AttributeTok{y       =} \ConstantTok{NULL}\NormalTok{,}
    \AttributeTok{fill    =} \StringTok{"Anomalía (°C)"}\NormalTok{,}
    \AttributeTok{caption =} \StringTok{"Fuente: NASA GISS"}
\NormalTok{  ) }\SpecialCharTok{+}
  \FunctionTok{theme\_minimal}\NormalTok{()}
\end{Highlighting}
\end{Shaded}

\begin{figure}[H]

{\centering \pandocbounded{\includegraphics[keepaspectratio]{visualizaciones_files/figure-pdf/unnamed-chunk-3-1.pdf}}

}

\caption{Mapa de calor de anomalías por mes}

\end{figure}%

\section{Gráfico 3 --- Comparación de
décadas}\label{gruxe1fico-3-comparaciuxf3n-de-duxe9cadas}

\begin{Shaded}
\begin{Highlighting}[]
\NormalTok{clima }\SpecialCharTok{|\textgreater{}}
  \FunctionTok{mutate}\NormalTok{(}\AttributeTok{decada =} \FunctionTok{paste0}\NormalTok{(}\FunctionTok{floor}\NormalTok{(year }\SpecialCharTok{/} \DecValTok{10}\NormalTok{) }\SpecialCharTok{*} \DecValTok{10}\NormalTok{, }\StringTok{"s"}\NormalTok{)) }\SpecialCharTok{|\textgreater{}}
  \FunctionTok{ggplot}\NormalTok{(}\FunctionTok{aes}\NormalTok{(}\AttributeTok{x =}\NormalTok{ anomalia, }\AttributeTok{fill =}\NormalTok{ decada)) }\SpecialCharTok{+}
  \FunctionTok{geom\_density}\NormalTok{(}\AttributeTok{alpha =} \FloatTok{0.5}\NormalTok{) }\SpecialCharTok{+}
  \FunctionTok{geom\_vline}\NormalTok{(}\AttributeTok{xintercept =} \DecValTok{0}\NormalTok{, }\AttributeTok{linetype =} \StringTok{"dashed"}\NormalTok{, }\AttributeTok{color =} \StringTok{"gray40"}\NormalTok{) }\SpecialCharTok{+}
  \FunctionTok{labs}\NormalTok{(}
    \AttributeTok{title    =} \StringTok{"Distribución de anomalías por década"}\NormalTok{,}
    \AttributeTok{subtitle =} \StringTok{"Cada curva representa una década"}\NormalTok{,}
    \AttributeTok{x        =} \StringTok{"Anomalía (°C)"}\NormalTok{,}
    \AttributeTok{y        =} \StringTok{"Densidad"}\NormalTok{,}
    \AttributeTok{fill     =} \StringTok{"Década"}\NormalTok{,}
    \AttributeTok{caption  =} \StringTok{"Fuente: NASA GISS"}
\NormalTok{  ) }\SpecialCharTok{+}
  \FunctionTok{theme\_minimal}\NormalTok{()}
\end{Highlighting}
\end{Shaded}

\begin{figure}[H]

{\centering \pandocbounded{\includegraphics[keepaspectratio]{visualizaciones_files/figure-pdf/unnamed-chunk-4-1.pdf}}

}

\caption{Distribución de anomalías por década}

\end{figure}%

\bookmarksetup{startatroot}

\chapter{Capítulo 3: Mapas y
Animaciones}\label{capuxedtulo-3-mapas-y-animaciones}

En este capítulo vamos a visualizar el cambio climático de forma
geográfica y animada usando \texttt{tmap} y \texttt{gganimate}.

\section{Cargar paquetes y datos}\label{cargar-paquetes-y-datos-1}

\begin{Shaded}
\begin{Highlighting}[]
\FunctionTok{library}\NormalTok{(tidyverse)}
\FunctionTok{library}\NormalTok{(lubridate)}
\FunctionTok{library}\NormalTok{(janitor)}
\FunctionTok{library}\NormalTok{(gganimate)}
\FunctionTok{library}\NormalTok{(gifski)}

\NormalTok{clima\_raw }\OtherTok{\textless{}{-}} \FunctionTok{read\_csv}\NormalTok{(}\StringTok{"GLB.Ts+dSST.csv"}\NormalTok{, }\AttributeTok{skip =} \DecValTok{1}\NormalTok{, }\AttributeTok{na =} \StringTok{"***"}\NormalTok{)}

\NormalTok{clima }\OtherTok{\textless{}{-}}\NormalTok{ clima\_raw }\SpecialCharTok{|\textgreater{}}
  \FunctionTok{clean\_names}\NormalTok{() }\SpecialCharTok{|\textgreater{}}
  \FunctionTok{select}\NormalTok{(year, jan}\SpecialCharTok{:}\NormalTok{dec) }\SpecialCharTok{|\textgreater{}}
  \FunctionTok{pivot\_longer}\NormalTok{(}\AttributeTok{cols =}\NormalTok{ jan}\SpecialCharTok{:}\NormalTok{dec,}
               \AttributeTok{names\_to  =} \StringTok{"mes"}\NormalTok{,}
               \AttributeTok{values\_to =} \StringTok{"anomalia"}\NormalTok{) }\SpecialCharTok{|\textgreater{}}
  \FunctionTok{mutate}\NormalTok{(}
    \AttributeTok{mes   =} \FunctionTok{match}\NormalTok{(mes, }\FunctionTok{c}\NormalTok{(}\StringTok{"jan"}\NormalTok{,}\StringTok{"feb"}\NormalTok{,}\StringTok{"mar"}\NormalTok{,}\StringTok{"apr"}\NormalTok{,}
                         \StringTok{"may"}\NormalTok{,}\StringTok{"jun"}\NormalTok{,}\StringTok{"jul"}\NormalTok{,}\StringTok{"aug"}\NormalTok{,}
                         \StringTok{"sep"}\NormalTok{,}\StringTok{"oct"}\NormalTok{,}\StringTok{"nov"}\NormalTok{,}\StringTok{"dec"}\NormalTok{)),}
    \AttributeTok{fecha =} \FunctionTok{make\_date}\NormalTok{(year, mes, }\DecValTok{1}\NormalTok{)}
\NormalTok{  ) }\SpecialCharTok{|\textgreater{}}
  \FunctionTok{drop\_na}\NormalTok{(anomalia)}
\end{Highlighting}
\end{Shaded}

\section{Anomalía promedio por
década}\label{anomaluxeda-promedio-por-duxe9cada}

\begin{Shaded}
\begin{Highlighting}[]
\NormalTok{clima\_decada }\OtherTok{\textless{}{-}}\NormalTok{ clima }\SpecialCharTok{|\textgreater{}}
  \FunctionTok{mutate}\NormalTok{(}\AttributeTok{decada =} \FunctionTok{floor}\NormalTok{(year }\SpecialCharTok{/} \DecValTok{10}\NormalTok{) }\SpecialCharTok{*} \DecValTok{10}\NormalTok{) }\SpecialCharTok{|\textgreater{}}
  \FunctionTok{group\_by}\NormalTok{(decada) }\SpecialCharTok{|\textgreater{}}
  \FunctionTok{summarise}\NormalTok{(}\AttributeTok{anomalia\_media =} \FunctionTok{mean}\NormalTok{(anomalia))}

\NormalTok{clima\_decada}
\end{Highlighting}
\end{Shaded}

\begin{verbatim}
# A tibble: 15 x 2
   decada anomalia_media
    <dbl>          <dbl>
 1   1880        -0.214 
 2   1890        -0.239 
 3   1900        -0.316 
 4   1910        -0.330 
 5   1920        -0.239 
 6   1930        -0.124 
 7   1940         0.0422
 8   1950        -0.048 
 9   1960        -0.0297
10   1970         0.0345
11   1980         0.246 
12   1990         0.385 
13   2000         0.587 
14   2010         0.805 
15   2020         1.07  
\end{verbatim}

\section{Animación --- evolución de la anomalía
anual}\label{animaciuxf3n-evoluciuxf3n-de-la-anomaluxeda-anual}

\begin{Shaded}
\begin{Highlighting}[]
\NormalTok{clima\_anual }\OtherTok{\textless{}{-}}\NormalTok{ clima }\SpecialCharTok{|\textgreater{}}
  \FunctionTok{group\_by}\NormalTok{(year) }\SpecialCharTok{|\textgreater{}}
  \FunctionTok{summarise}\NormalTok{(}\AttributeTok{anomalia\_anual =} \FunctionTok{mean}\NormalTok{(anomalia))}

\NormalTok{animacion }\OtherTok{\textless{}{-}}\NormalTok{ clima\_anual }\SpecialCharTok{|\textgreater{}}
  \FunctionTok{ggplot}\NormalTok{(}\FunctionTok{aes}\NormalTok{(}\AttributeTok{x =}\NormalTok{ year, }\AttributeTok{y =}\NormalTok{ anomalia\_anual)) }\SpecialCharTok{+}
  \FunctionTok{geom\_line}\NormalTok{(}\AttributeTok{color =} \StringTok{"\#2196F3"}\NormalTok{, }\AttributeTok{linewidth =} \FloatTok{0.8}\NormalTok{) }\SpecialCharTok{+}
  \FunctionTok{geom\_point}\NormalTok{(}\AttributeTok{color =} \StringTok{"\#E53935"}\NormalTok{, }\AttributeTok{size =} \DecValTok{2}\NormalTok{) }\SpecialCharTok{+}
  \FunctionTok{geom\_hline}\NormalTok{(}\AttributeTok{yintercept =} \DecValTok{0}\NormalTok{, }\AttributeTok{linetype =} \StringTok{"dashed"}\NormalTok{, }\AttributeTok{color =} \StringTok{"gray50"}\NormalTok{) }\SpecialCharTok{+}
  \FunctionTok{labs}\NormalTok{(}
    \AttributeTok{title   =} \StringTok{"Anomalía de temperatura global — Año: \{frame\_along\}"}\NormalTok{,}
    \AttributeTok{x       =} \StringTok{"Año"}\NormalTok{,}
    \AttributeTok{y       =} \StringTok{"Anomalía (°C)"}\NormalTok{,}
    \AttributeTok{caption =} \StringTok{"Fuente: NASA GISS"}
\NormalTok{  ) }\SpecialCharTok{+}
  \FunctionTok{theme\_minimal}\NormalTok{() }\SpecialCharTok{+}
  \FunctionTok{transition\_reveal}\NormalTok{(year)}

\FunctionTok{animate}\NormalTok{(animacion,}
        \AttributeTok{nframes  =} \DecValTok{100}\NormalTok{,}
        \AttributeTok{fps      =} \DecValTok{10}\NormalTok{,}
        \AttributeTok{width    =} \DecValTok{800}\NormalTok{,}
        \AttributeTok{height   =} \DecValTok{400}\NormalTok{,}
        \AttributeTok{units    =} \StringTok{"px"}\NormalTok{,}
        \AttributeTok{res      =} \DecValTok{96}\NormalTok{,}
        \AttributeTok{renderer =} \FunctionTok{gifski\_renderer}\NormalTok{())}
\end{Highlighting}
\end{Shaded}

\pandocbounded{\includegraphics[keepaspectratio]{mapas_files/figure-pdf/unnamed-chunk-3-1.gif}}

\section{Gráfico de barras por
década}\label{gruxe1fico-de-barras-por-duxe9cada}

\begin{Shaded}
\begin{Highlighting}[]
\NormalTok{clima\_decada }\SpecialCharTok{|\textgreater{}}
  \FunctionTok{ggplot}\NormalTok{(}\FunctionTok{aes}\NormalTok{(}\AttributeTok{x =} \FunctionTok{factor}\NormalTok{(decada),}
             \AttributeTok{y =}\NormalTok{ anomalia\_media,}
             \AttributeTok{fill =}\NormalTok{ anomalia\_media)) }\SpecialCharTok{+}
  \FunctionTok{geom\_col}\NormalTok{() }\SpecialCharTok{+}
  \FunctionTok{scale\_fill\_gradient2}\NormalTok{(}
    \AttributeTok{low      =} \StringTok{"\#2166AC"}\NormalTok{,}
    \AttributeTok{mid      =} \StringTok{"white"}\NormalTok{,}
    \AttributeTok{high     =} \StringTok{"\#D6604D"}\NormalTok{,}
    \AttributeTok{midpoint =} \DecValTok{0}
\NormalTok{  ) }\SpecialCharTok{+}
  \FunctionTok{labs}\NormalTok{(}
    \AttributeTok{title    =} \StringTok{"Anomalía promedio por década"}\NormalTok{,}
    \AttributeTok{subtitle =} \StringTok{"El calentamiento se acelera desde 1980"}\NormalTok{,}
    \AttributeTok{x        =} \StringTok{"Década"}\NormalTok{,}
    \AttributeTok{y        =} \StringTok{"Anomalía media (°C)"}\NormalTok{,}
    \AttributeTok{fill     =} \StringTok{"Anomalía"}\NormalTok{,}
    \AttributeTok{caption  =} \StringTok{"Fuente: NASA GISS"}
\NormalTok{  ) }\SpecialCharTok{+}
  \FunctionTok{theme\_minimal}\NormalTok{() }\SpecialCharTok{+}
  \FunctionTok{theme}\NormalTok{(}\AttributeTok{legend.position =} \StringTok{"none"}\NormalTok{)}
\end{Highlighting}
\end{Shaded}

\pandocbounded{\includegraphics[keepaspectratio]{mapas_files/figure-pdf/unnamed-chunk-4-1.png}}

\bookmarksetup{startatroot}

\chapter{Capítulo 4: Series de Tiempo y
Proyecciones}\label{capuxedtulo-4-series-de-tiempo-y-proyecciones}

En este capítulo vamos a analizar la tendencia de la temperatura global
usando series de tiempo y proyectar cómo podría evolucionar en los
próximos 20 años.

\section{Cargar paquetes y datos}\label{cargar-paquetes-y-datos-2}

\begin{Shaded}
\begin{Highlighting}[]
\FunctionTok{library}\NormalTok{(tidyverse)}
\FunctionTok{library}\NormalTok{(lubridate)}
\FunctionTok{library}\NormalTok{(janitor)}
\FunctionTok{library}\NormalTok{(tsibble)}
\FunctionTok{library}\NormalTok{(fable)}
\FunctionTok{library}\NormalTok{(feasts)}

\NormalTok{clima\_raw }\OtherTok{\textless{}{-}} \FunctionTok{read\_csv}\NormalTok{(}\StringTok{"GLB.Ts+dSST.csv"}\NormalTok{, }\AttributeTok{skip =} \DecValTok{1}\NormalTok{, }\AttributeTok{na =} \StringTok{"***"}\NormalTok{)}

\NormalTok{clima }\OtherTok{\textless{}{-}}\NormalTok{ clima\_raw }\SpecialCharTok{|\textgreater{}}
  \FunctionTok{clean\_names}\NormalTok{() }\SpecialCharTok{|\textgreater{}}
  \FunctionTok{select}\NormalTok{(year, jan}\SpecialCharTok{:}\NormalTok{dec) }\SpecialCharTok{|\textgreater{}}
  \FunctionTok{pivot\_longer}\NormalTok{(}\AttributeTok{cols =}\NormalTok{ jan}\SpecialCharTok{:}\NormalTok{dec,}
               \AttributeTok{names\_to  =} \StringTok{"mes"}\NormalTok{,}
               \AttributeTok{values\_to =} \StringTok{"anomalia"}\NormalTok{) }\SpecialCharTok{|\textgreater{}}
  \FunctionTok{mutate}\NormalTok{(}
    \AttributeTok{mes   =} \FunctionTok{match}\NormalTok{(mes, }\FunctionTok{c}\NormalTok{(}\StringTok{"jan"}\NormalTok{,}\StringTok{"feb"}\NormalTok{,}\StringTok{"mar"}\NormalTok{,}\StringTok{"apr"}\NormalTok{,}
                         \StringTok{"may"}\NormalTok{,}\StringTok{"jun"}\NormalTok{,}\StringTok{"jul"}\NormalTok{,}\StringTok{"aug"}\NormalTok{,}
                         \StringTok{"sep"}\NormalTok{,}\StringTok{"oct"}\NormalTok{,}\StringTok{"nov"}\NormalTok{,}\StringTok{"dec"}\NormalTok{)),}
    \AttributeTok{fecha =} \FunctionTok{make\_date}\NormalTok{(year, mes, }\DecValTok{1}\NormalTok{)}
\NormalTok{  ) }\SpecialCharTok{|\textgreater{}}
  \FunctionTok{drop\_na}\NormalTok{(anomalia)}
\end{Highlighting}
\end{Shaded}

\section{Convertir a serie de tiempo}\label{convertir-a-serie-de-tiempo}

\begin{Shaded}
\begin{Highlighting}[]
\NormalTok{clima\_ts }\OtherTok{\textless{}{-}}\NormalTok{ clima }\SpecialCharTok{|\textgreater{}}
  \FunctionTok{mutate}\NormalTok{(mes\_año }\OtherTok{=} \FunctionTok{yearmonth}\NormalTok{(fecha)) }\SpecialCharTok{|\textgreater{}}
  \FunctionTok{as\_tsibble}\NormalTok{(}\AttributeTok{index =}\NormalTok{ mes\_año)}

\NormalTok{clima\_ts}
\end{Highlighting}
\end{Shaded}

\begin{verbatim}
# A tsibble: 1,756 x 5 [1M]
    year   mes anomalia fecha        mes_año
   <dbl> <int>    <dbl> <date>         <mth>
 1  1880     1    -0.19 1880-01-01 1880 ene.
 2  1880     2    -0.25 1880-02-01 1880 feb.
 3  1880     3    -0.1  1880-03-01 1880 mar.
 4  1880     4    -0.17 1880-04-01 1880 abr.
 5  1880     5    -0.1  1880-05-01 1880 may.
 6  1880     6    -0.21 1880-06-01 1880 jun.
 7  1880     7    -0.19 1880-07-01 1880 jul.
 8  1880     8    -0.11 1880-08-01 1880 ago.
 9  1880     9    -0.15 1880-09-01 1880 sep.
10  1880    10    -0.24 1880-10-01 1880 oct.
# i 1,746 more rows
\end{verbatim}

\section{Explorar la serie de tiempo}\label{explorar-la-serie-de-tiempo}

\begin{Shaded}
\begin{Highlighting}[]
\NormalTok{clima\_ts }\SpecialCharTok{|\textgreater{}}
  \FunctionTok{autoplot}\NormalTok{(anomalia) }\SpecialCharTok{+}
  \FunctionTok{labs}\NormalTok{(}
    \AttributeTok{title   =} \StringTok{"Anomalía de temperatura global mensual"}\NormalTok{,}
    \AttributeTok{x       =} \StringTok{"Fecha"}\NormalTok{,}
    \AttributeTok{y       =} \StringTok{"Anomalía (°C)"}\NormalTok{,}
    \AttributeTok{caption =} \StringTok{"Fuente: NASA GISS"}
\NormalTok{  ) }\SpecialCharTok{+}
  \FunctionTok{theme\_minimal}\NormalTok{()}
\end{Highlighting}
\end{Shaded}

\pandocbounded{\includegraphics[keepaspectratio]{proyecciones_files/figure-pdf/unnamed-chunk-3-1.pdf}}

\section{Descomposición de la
serie}\label{descomposiciuxf3n-de-la-serie}

\begin{Shaded}
\begin{Highlighting}[]
\NormalTok{clima\_ts }\SpecialCharTok{|\textgreater{}}
  \FunctionTok{model}\NormalTok{(}\FunctionTok{STL}\NormalTok{(anomalia)) }\SpecialCharTok{|\textgreater{}}
  \FunctionTok{components}\NormalTok{() }\SpecialCharTok{|\textgreater{}}
  \FunctionTok{autoplot}\NormalTok{() }\SpecialCharTok{+}
  \FunctionTok{labs}\NormalTok{(}
    \AttributeTok{title   =} \StringTok{"Descomposición STL de la anomalía"}\NormalTok{,}
    \AttributeTok{caption =} \StringTok{"Tendencia + Estacionalidad + Residuo"}
\NormalTok{  ) }\SpecialCharTok{+}
  \FunctionTok{theme\_minimal}\NormalTok{()}
\end{Highlighting}
\end{Shaded}

\pandocbounded{\includegraphics[keepaspectratio]{proyecciones_files/figure-pdf/unnamed-chunk-4-1.pdf}}

\section{Proyección a 20 años}\label{proyecciuxf3n-a-20-auxf1os}

\begin{Shaded}
\begin{Highlighting}[]
\NormalTok{modelo }\OtherTok{\textless{}{-}}\NormalTok{ clima\_ts }\SpecialCharTok{|\textgreater{}}
  \FunctionTok{model}\NormalTok{(}
    \AttributeTok{arima =} \FunctionTok{ARIMA}\NormalTok{(anomalia),}
    \AttributeTok{naive =} \FunctionTok{NAIVE}\NormalTok{(anomalia)}
\NormalTok{  )}

\NormalTok{proyeccion }\OtherTok{\textless{}{-}}\NormalTok{ modelo }\SpecialCharTok{|\textgreater{}}
  \FunctionTok{forecast}\NormalTok{(}\AttributeTok{h =} \StringTok{"20 years"}\NormalTok{)}

\NormalTok{proyeccion }\SpecialCharTok{|\textgreater{}}
  \FunctionTok{filter}\NormalTok{(.model }\SpecialCharTok{==} \StringTok{"arima"}\NormalTok{) }\SpecialCharTok{|\textgreater{}}
  \FunctionTok{autoplot}\NormalTok{(clima\_ts) }\SpecialCharTok{+}
  \FunctionTok{labs}\NormalTok{(}
    \AttributeTok{title    =} \StringTok{"Proyección de anomalía de temperatura (próximos 20 años)"}\NormalTok{,}
    \AttributeTok{subtitle =} \StringTok{"Modelo ARIMA · Banda sombreada = intervalo de confianza"}\NormalTok{,}
    \AttributeTok{x        =} \StringTok{"Fecha"}\NormalTok{,}
    \AttributeTok{y        =} \StringTok{"Anomalía (°C)"}\NormalTok{,}
    \AttributeTok{caption  =} \StringTok{"Fuente: NASA GISS · Proyección con fable"}
\NormalTok{  ) }\SpecialCharTok{+}
  \FunctionTok{theme\_minimal}\NormalTok{()}
\end{Highlighting}
\end{Shaded}

\pandocbounded{\includegraphics[keepaspectratio]{proyecciones_files/figure-pdf/unnamed-chunk-5-1.pdf}}

\bookmarksetup{startatroot}

\chapter{Conclusiones}\label{conclusiones}

A lo largo de este book analizamos la evolución de la temperatura global
desde 1880 hasta hoy usando datos reales de la NASA y herramientas de R.

\section{Lo que encontramos}\label{lo-que-encontramos}

Los datos muestran de forma clara y contundente:

\begin{itemize}
\tightlist
\item
  La temperatura global aumentó aproximadamente \textbf{1.2°C} desde
  1880
\item
  El calentamiento se \textbf{aceleró notablemente desde 1980}
\item
  Los últimos años registran las anomalías \textbf{más altas de toda la
  serie}
\item
  La tendencia proyectada indica que el calentamiento
  \textbf{continuará}
\end{itemize}

\section{Lo que aprendimos en R}\label{lo-que-aprendimos-en-r}

{\def\LTcaptype{none} % do not increment counter
\begin{longtable}[]{@{}
  >{\raggedright\arraybackslash}p{(\linewidth - 4\tabcolsep) * \real{0.1299}}
  >{\raggedright\arraybackslash}p{(\linewidth - 4\tabcolsep) * \real{0.4286}}
  >{\raggedright\arraybackslash}p{(\linewidth - 4\tabcolsep) * \real{0.4416}}@{}}
\toprule\noalign{}
\begin{minipage}[b]{\linewidth}\raggedright
Capítulo
\end{minipage} & \begin{minipage}[b]{\linewidth}\raggedright
Habilidad
\end{minipage} & \begin{minipage}[b]{\linewidth}\raggedright
Librería
\end{minipage} \\
\midrule\noalign{}
\endhead
\bottomrule\noalign{}
\endlastfoot
1 & Importar y limpiar datos reales & \texttt{readr}, \texttt{janitor},
\texttt{dplyr} \\
2 & Visualizaciones estáticas & \texttt{ggplot2}, \texttt{scales},
\texttt{patchwork} \\
3 & Animaciones & \texttt{gganimate}, \texttt{gifski} \\
4 & Series de tiempo y proyecciones & \texttt{tsibble}, \texttt{fable},
\texttt{feasts} \\
\end{longtable}
}

\section{Fuentes}\label{fuentes}

\begin{itemize}
\tightlist
\item
  \href{https://data.giss.nasa.gov/gistemp/}{NASA GISS Surface
  Temperature Analysis}
\item
  \href{https://ourworldindata.org/climate-change}{Our World in Data ---
  Climate Change}
\item
  Wickham, H. (2016). \emph{ggplot2: Elegant Graphics for Data Analysis}
\end{itemize}




\end{document}
